\chapter{Dynamic Hosting Capacity and Spectral Planning Levers}

The preceding chapters define the damping operator hierarchy. This chapter
uses that hierarchy to formulate the planning questions that motivated the
project: how much inverter-based resource (IBR) conversion can be hosted, where
virtual inertia or virtual damping should be placed, and which network
reinforcements can improve robustness before resorting to batteries or other
large storage devices. The common theme is that every planning answer must be
attached to a margin. Voltage feasibility, graph connectivity, damping ratio,
protected resolvent gain, and controller-resonance distance are different
objects.

\section{From Static Hosting to Spectral Hosting}

Let an operating point be described by complex bus voltages
\(V_i=|V_i|e^{j\theta_i}\) and bus admittance matrix
\(Y_{\rm bus}=G+jB\). The exact AC active- and reactive-power injections are
\begin{align}
  P_i(V,\theta)
  &=
  |V_i|\sum_{j=1}^n |V_j|
  \left(
  G_{ij}\cos(\theta_i-\theta_j)
  +
  B_{ij}\sin(\theta_i-\theta_j)
  \right),
  \label{eq:ac_active_power}\\
  Q_i(V,\theta)
  &=
  |V_i|\sum_{j=1}^n |V_j|
  \left(
  G_{ij}\sin(\theta_i-\theta_j)
  -
  B_{ij}\cos(\theta_i-\theta_j)
  \right).
  \label{eq:ac_reactive_power}
\end{align}
The linearized active-power stiffness is the angle block of the AC Jacobian.
After removing the uniform angle reference, write the corresponding positive
semidefinite stiffness as \(L_{\rm op}\). In the lossless small-angle
approximation, \(L_{\rm op}\) reduces to a weighted Laplacian.

A static hosting set for a vector of converter injections or conversions
\(\rho\in\mathcal U\) can be written abstractly as
\begin{equation}
  \mathcal H_{\rm stat}
  =
  \{\rho\in\mathcal U:\ c_{\rm V}(\rho)\le0,\ 
  c_{\rm th}(\rho)\le0,\ c_{\rm q}(\rho)\le0\},
  \label{eq:static_hosting_set}
\end{equation}
where the constraints may represent voltage, thermal, protection, power
quality, or operational limits. This is the classical hosting-capacity
language. It is indispensable, but it does not by itself certify small-signal
damping.

\begin{definition}[Dynamic spectral hosting set]
For a target damping ratio \(\zetaReq\), protected resolvent threshold
\(g_{\rm req}\), and admissible static set \(\mathcal H_{\rm stat}\), define
the dynamic spectral hosting set
\begin{equation}
  \mathcal H_{\rm dyn}
  =
  \left\{
  \rho\in\mathcal H_{\rm stat}:
  \zeta_{\min}(\rho)\ge \zetaReq,\ 
  g_{\rm prot}(\rho)\le g_{\rm req},\
  T(s;\rho)\hbox{ has no forbidden roots}
  \right\}.
  \label{eq:dynamic_hosting_set}
\end{equation}
For a penetration direction \(\alpha\ge0\), the directional dynamic hosting
capacity is
\begin{equation}
  h_{\rm dyn}(\alpha)
  =
  \sup\{\eta\ge0:\ \eta\alpha\in\mathcal H_{\rm dyn}\}.
  \label{eq:directional_dynamic_hosting}
\end{equation}
\end{definition}

This definition deliberately extends time-varying hosting-capacity ideas
\cite{oliveira2019dynamicHosting,jain2020dynamicHosting}. In distribution
planning, dynamic hosting often means that voltage or thermal violations are
counted with their duration and severity. Here it means that the electromechanical
and controller spectra also impose capacity limits. A system can be voltage
safe and still dynamically unsafe if the critical pole violates
\(\zetaReq\), if the protected resolvent gain is too large, or if a stable
network-family pole lies close to a controller pole with large self-energy
sensitivity.

\section{First-Order Spectral Hosting Screen}

In the fixed-\(D\) layer, let
\begin{equation}
  L(\rho)\phi_k(\rho)
  =
  \nu_k(\rho)M(\rho)\phi_k(\rho),
  \qquad
  \phi_k(\rho)^TM(\rho)\phi_\ell(\rho)=\delta_{k\ell}.
  \label{eq:hosting_generalized_modes}
\end{equation}
The diagonal damping-region screen from Chapter~4 is
\begin{equation}
  h_k(\rho)
  =
  \delta_k(\rho)-2\zetaReq\sqrt{\nu_k(\rho)}
  \ge0,
  \qquad
  \delta_k(\rho)=\phi_k(\rho)^TD(\rho)\phi_k(\rho).
  \label{eq:hosting_modal_gap}
\end{equation}
The quantity \(h_k\) is a modal headroom, not an eigenvalue. It is exact only
when the relevant modes decouple and becomes a screen otherwise.

For a simple generalized mode and small perturbation \(\rho_a\),
\begin{equation}
  \frac{\partial \nu_k}{\partial \rho_a}
  =
  \phi_k^T
  \left(
  L_a-\nu_k M_a
  \right)
  \phi_k,
  \label{eq:hosting_nu_derivative}
\end{equation}
where \(L_a=\partial L/\partial\rho_a\) and
\(M_a=\partial M/\partial\rho_a\). The direct damping derivative is
\(\phi_k^TD_a\phi_k\), with modal-rotation corrections when damping is
non-proportional. The first-order hosting-gap derivative is therefore
\begin{equation}
  g_{ka}
  =
  \frac{\partial h_k}{\partial\rho_a}
  =
  \frac{\partial\delta_k}{\partial\rho_a}
  -
  \frac{\zetaReq}{\sqrt{\nu_k}}
  \frac{\partial\nu_k}{\partial\rho_a}.
  \label{eq:hosting_gap_gradient}
\end{equation}

\begin{proposition}[Linear spectral hosting screen]
Assume \(h_k(0)>0\) for every monitored mode \(k\in\mathcal K\), and assume
that the first-order expansion of \eqref{eq:hosting_modal_gap} is used as a
tangent screen. Define
\begin{equation}
  A_{ka}=-g_{ka},
  \qquad
  b_k=h_k(0).
  \label{eq:hosting_polytope_matrices}
\end{equation}
Then the polytope
\begin{equation}
  \mathcal P_{\rm host}^{(1)}
  =
  \{\rho\in\mathcal U:\ A\rho\le b\}
  \label{eq:hosting_polytope}
\end{equation}
is the first-order spectral hosting screen associated with the damping-region
margin. It is not, by itself, a certified inner approximation of
\(\mathcal H_{\rm dyn}\).
\end{proposition}

\begin{proof}
The first-order expansion gives
\[
  h_k(\rho)=h_k(0)+\sum_a g_{ka}\rho_a+O(\norm{\rho}^2).
\]
Dropping the higher-order term and imposing \(h_k(\rho)\ge0\) gives
\(-\sum_ag_{ka}\rho_a\le h_k(0)\), which is \eqref{eq:hosting_polytope}.
\end{proof}

\begin{proposition}[Curvature-certified hosting region]
Suppose that, on a trust region \(\mathcal B_r=\{\rho:\norm{\rho}_2\le r\}\),
each monitored headroom satisfies
\[
  \norm{\nabla^2 h_k(\rho)}_2\le \beta_k .
\]
Then every \(\rho\in\mathcal U\cap\mathcal B_r\) satisfying
\begin{equation}
  h_k(0)+g_k^T\rho-\frac{\beta_k}{2}\norm{\rho}_2^2\ge0,
  \qquad k\in\mathcal K,
  \label{eq:hosting_curvature_certificate}
\end{equation}
also satisfies the second-order lower bound for the monitored modal headrooms.
\end{proposition}

\begin{proof}
Taylor's theorem with the stated Hessian bound gives
\[
  h_k(\rho)
  \ge
  h_k(0)+g_k^T\rho-\frac{\beta_k}{2}\norm{\rho}_2^2
\]
for all \(\rho\in\mathcal B_r\). Imposing
\eqref{eq:hosting_curvature_certificate} therefore certifies nonnegative
headroom for the monitored screen.
\end{proof}

The certified set is generally quadratic or conic, not a polytope. This is the
right object when the claim is an inner dynamic-hosting region rather than a
first-order ranking.

\section{Certified Hosting and Model Escalation}

A dynamic-hosting screen becomes a certified hosting statement only after the
model level is certified over the protected region. Let \(\Gamma\) enclose the
set of eigenvalues whose damping ratio, abscissa, or controller-resonance
distance limits hosting. For a reduced screen \(\widehat T(s;\rho)\), define
\begin{equation}
  \eta_\Gamma(\rho)
  =
  \sup_{s\in\Gamma}
  \left\|
  \widehat T(s;\rho)^{-1}
  \left[T(s;\rho)-\widehat T(s;\rho)\right]
  \right\|_2 .
  \label{eq:hosting_eta_gamma}
\end{equation}
If \(\eta_\Gamma(\rho)<1\), the reduced model preserves the protected
eigenvalue count inside \(\Gamma\). If \(\eta_\Gamma(\rho)\ge1\), the hosting
decision must escalate from the diagonal screen to the QEP, from the QEP to the
Schur NEP, or from the identified NEP to the full DAE, depending on where the
uncertainty enters.

For finite penetration increments, the same principle should be stated as a
fragility budget. Given a protected boundary \(\partial\mathcal S\), define
\begin{equation}
  r_{\rm host}(\rho)
  =
  \inf
  \left\{
  \|W_u u\|:
  \exists s\in\partial\mathcal S,\ 
  \det T(s;\rho,u)=0
  \right\}.
  \label{eq:hosting_fragility_budget}
\end{equation}
A hosting point is not robust merely because the first-order headroom is
positive. It is robust to the declared uncertainty family only if
\(r_{\rm host}\) exceeds the uncertainty or action budget assigned to that
point. This is the dynamic analogue of a voltage or thermal security margin.

\begin{remark}[SG-to-IBR conversion]
If action \(\rho_i\) replaces synchronous inertia and damping at bus \(i\)
without changing the network stiffness, then locally
\[
  M_i(\rho_i)=M_i^0-\rho_i\Delta M_i,\qquad
  D_i(\rho_i)=D_i^0-\rho_i\Delta D_i.
\]
The direct mode-frequency contribution is
\[
  \frac{\partial\nu_k}{\partial\rho_i}
  =
  \nu_k\Delta M_i\,\phi_k(i)^2.
\]
The direct damping contribution is approximately
\[
  \frac{\partial\delta_k}{\partial\rho_i}
  \approx
  -\Delta D_i\,\phi_k(i)^2
  +
  \hbox{modal-rotation terms}.
\]
Therefore SG-to-IBR conversion is most restrictive where the monitored mode
has large \(\phi_k(i)^2\) and the removed damping-to-inertia package is large.
This is the mathematical reason a hosting limit is not only a voltage limit.
\end{remark}

\section{Where to Place Virtual Inertia}

Virtual inertia is not a single planning lever. It may mean inertia-only
emulation, virtual synchronous-machine control with damping, grid-forming
droop with filtering, or a controller whose reduced effect is frequency
dependent. The fixed-\(D\) screen is still useful because it identifies what
would be true if the virtual device were represented locally by added mass and
added damping:
\begin{equation}
  M_i(\mu_i)=M_i^0+\mu_i,\qquad
  D_i(\mu_i)=D_i^0+\gamma_i^v\mu_i.
  \label{eq:virtual_inertia_local_model}
\end{equation}
Using the fixed-mode approximations of Chapter~4,
\begin{equation}
  \frac{\partial h_k}{\partial \mu_i}
  \approx
  \left(\gamma_i^v-\delta_k+\zetaReq\sqrt{\nu_k}\right)\phi_k(i)^2.
  \label{eq:virtual_inertia_hosting_gradient}
\end{equation}
Thus inertia-only support \((\gamma_i^v=0)\) can reduce the damping-region
headroom near the boundary \( \delta_k\approx2\zetaReq\sqrt{\nu_k}\). The same
bus can still be excellent for RoCoF or frequency-nadir objectives because
inertia reduces modal acceleration. It is not automatically excellent for
damping ratio.

For a set of monitored modes, a placement score that respects this distinction
is
\begin{equation}
  J_i^{\rm VI}
  =
  \sum_{k\in\mathcal K}\omega_k
  \left[
  \eta_{\rm f}\nu_k\phi_k(i)^2
  +
  \eta_{\rm h}
  \left(\gamma_i^v-\delta_k+\zetaReq\sqrt{\nu_k}\right)\phi_k(i)^2
  \right],
  \label{eq:virtual_inertia_score}
\end{equation}
where the first term rewards frequency-security leverage and the second term
rewards damping-headroom leverage. The weights \(\eta_{\rm f}\) and
\(\eta_{\rm h}\) must be declared before ranking. Otherwise the ranking hides
a tradeoff between slower frequency motion and better modal damping.

\section{Line-Shunt Taxonomy and Substitutability}

The project also asks whether the network can be improved before installing
batteries or local inertial devices. A structural answer comes from separating
line and shunt directions.

Let \(b_{ij}=e_i-e_j\) and define
\begin{equation}
  \mathcal L_{\rm all}
  =
  \operatorname{span}\{b_{ij}b_{ij}^T:\ 1\le i<j\le n\},
  \qquad
  \mathcal D_{\rm sh}
  =
  \{\operatorname{diag}(g):\ g\in\R^n\}.
  \label{eq:line_shunt_spaces}
\end{equation}

\begin{theorem}[Line-shunt decomposition]
Every real symmetric matrix \(X\in\R^{n\times n}\) has a unique decomposition
\begin{equation}
  X=X_{\mathcal L}+X_{\mathcal D},
  \qquad
  X_{\mathcal L}\in\mathcal L_{\rm all},\quad
  X_{\mathcal D}\in\mathcal D_{\rm sh}.
  \label{eq:line_shunt_decomposition}
\end{equation}
Equivalently,
\[
  \operatorname{Sym}(n)=\mathcal L_{\rm all}\oplus\mathcal D_{\rm sh}.
\]
\end{theorem}

\begin{proof}
Set the off-diagonal entries of \(X_{\mathcal L}\) equal to those of \(X\) and
choose the diagonal of \(X_{\mathcal L}\) so that every row sums to zero. Then
\(X_{\mathcal L}\) is a Laplacian-type matrix and
\(X_{\mathcal D}=X-X_{\mathcal L}\) is diagonal. The intersection of the two
spaces is zero because a diagonal zero-row-sum matrix must be zero.
\end{proof}

For local inertia,
\begin{equation}
  \frac{\partial\Lt}{\partial M_i}
  =
  -\frac{1}{2M_i}
  \left(E_i\Lt+\Lt E_i\right),
  \qquad
  E_i=e_ie_i^T.
  \label{eq:local_inertia_operator_location}
\end{equation}
This perturbation cannot create off-diagonal support on a non-existing edge.
Therefore it lies structurally in existing-line directions plus a shunt
residual:
\begin{equation}
  \frac{\partial\Lt}{\partial M_i}
  \in
  \mathcal L_{\rm ex}\oplus\mathcal D_{\rm sh}.
  \label{eq:inertia_existing_shunt}
\end{equation}
It has no new-line component. This statement does not mean inertia and line
reinforcement are dynamically equivalent. It says where the inertia action
lives in the perturbation space.

Let
\begin{equation}
  -\frac{\partial\Lt}{\partial M_i}
  =
  X_i^{\rm ex}+\operatorname{diag}(r_i),
  \qquad
  X_i^{\rm ex}\in\mathcal L_{\rm ex}.
  \label{eq:inertia_shunt_residual}
\end{equation}
For a critical mode \(q_c\), define the modal substitutability score
\begin{equation}
  s_i
  =
  1-
  \frac{|r_i^T(q_c\odot q_c)|}
  {\norm{r_i}_2\,\norm{q_c\odot q_c}_2+\varepsilon_{\rm reg}}.
  \label{eq:modal_substitutability_score}
\end{equation}
High \(s_i\) means the shunt residual is nearly invisible to the critical
mode; existing-line actions may be tried before local inertia. Low \(s_i\)
means the shunt residual aligns with the critical mode; topology-only actions
are less likely to substitute for a local inertial or synchronous-condenser
lever. Internal random-graph checks in \texttt{temp/poster/} verified the
algebraic residual of this decomposition to approximately \(10^{-14}\), which
is a structural verification rather than a power-system performance proof.
The score is still not a feasibility certificate for construction. A physical
line-only substitute must solve a constrained projection with nonnegative line
increments, corridor limits, and costs; membership in a linear span is only the
algebraic first screen.
\begin{equation}
  \min_{\Delta b\in\mathcal C_{\rm line}}
  \left\|
  -\frac{\partial\Lt}{\partial M_i}
  -
  \sum_{e\in\mathcal E_{\rm cand}}\Delta b_e a_ea_e^T
  \right\|_W,
  \qquad
  \mathcal C_{\rm line}=\{\Delta b\ge0:\ c^T\Delta b\le B,\ 
  \Delta b_e\le \bar b_e\}.
  \label{eq:line_inertia_constrained_projection}
\end{equation}

\section{Reinforcing Lines Without Batteries}

For a line \(e=(i,j)\), the modal-frequency derivative is
\begin{equation}
  \frac{\partial\nu_k}{\partial b_e}
  =
  (\phi_k(i)-\phi_k(j))^2.
  \label{eq:hosting_line_frequency}
\end{equation}
The damping-headroom derivative is
\begin{equation}
  \frac{\partial h_k}{\partial b_e}
  =
  \frac{\partial\delta_k}{\partial b_e}
  -
  \frac{\zetaReq}{\sqrt{\nu_k}}
  (\phi_k(i)-\phi_k(j))^2.
  \label{eq:hosting_line_headroom}
\end{equation}
This equation is the cleanest answer to the line-reinforcement question. A
line that strongly raises modal stiffness can still reduce damping headroom
unless it also raises modal damping through eigenvector rotation, damping
placement, or controller action.

The network-only benefits should still be used. If the objective is to improve
algebraic connectivity, use \(G_e^{(2)}=(\phi_2(i)-\phi_2(j))^2\). If the
objective is to reduce total effective resistance, use
\begin{equation}
  G_e^{(K)}
  =
  a_e^T(L^+)^2a_e,
  \qquad
  \frac{\partial}{\partial b_e}
  n\operatorname{trace}(L^+)
  =
  -nG_e^{(K)}.
  \label{eq:effective_resistance_line_gain}
\end{equation}
But a dynamic-hosting line screen should report the vector
\begin{equation}
  \ell_e
  =
  \begin{bmatrix}
  G_e^{(2)}\\
  G_e^{(K)}\\
  -\min_{k\in\mathcal K}\partial h_k/\partial b_e\\
  \partial S_g/\partial b_e\\
  F_e
  \end{bmatrix},
  \label{eq:line_screen_vector}
\end{equation}
where \(F_e\) is a feasibility penalty. A line should be called a robustness
improvement only after the selected margin improves and feasibility is
acceptable. Otherwise it is only a graph-stiffness improvement.

For controller-aware planning, the line report should be extended to a safe
reinforcement envelope:
\begin{equation}
  \mathcal R_e
  =
  \begin{bmatrix}
  G_e^{(2)}\\
  G_e^{(K)}\\
  \zeta_{p,e}^{\rm static}\\
  \zeta_{p,e}^{\rm NEP}\\
  \rho_{{\rm ctrl},e}\\
  \rho_{{\rm mem},e}\\
  p_{e,\max}^{\rm safe}
  \end{bmatrix}.
  \label{eq:line_safe_reinforcement_envelope}
\end{equation}
Here \(\rho_{{\rm ctrl},e}\) is a magnitude-normalized controller-channel
dominance score, \(\rho_{{\rm mem},e}\) is the memory-dominance score induced
by \(\Pi_s(s)\), and \(p_{e,\max}^{\rm safe}\) is the local action budget before
a protected margin is exhausted. The controller-distance term is
\(d^{\rm ctrl}_{k\ell}=|s_k-\mu_\ell|/(|s_k|+|\mu_\ell|)\), where
\(\mu_\ell\) is a condensed controller pole. A first-order version is
\begin{equation}
  p_{e,\max}^{\rm safe}
  =
  \min_{k,\ell}
  \left\{
  \frac{\zeta_k-\zetaReq}
       {[-\zeta_{p,e,k}^{\rm NEP}]_+},
  \frac{d^{\rm ctrl}_{k\ell}-d_{\rm req}}
       {[-d_{p,e,k\ell}^{\rm ctrl}]_+}
  \right\},
  \label{eq:line_safe_reinforcement_budget}
\end{equation}
with the convention that a zero denominator gives no active local restriction.
Finite upgrades must still be checked by continuation or full eigenvalue
recomputation.

This envelope separates four planning classes:
\begin{enumerate}
\item Safe reinforcement: graph strength and protected damping both improve.
\item Ordinary tradeoff: graph strength improves, but damping ratio falls for
the static stiffness-frequency reason.
\item Control-amplified reversal: the controller channel worsens a degradation
already visible in the static screen.
\item Controller-induced reversal: the static screen accepts the line but the
controller-aware model rejects it.
\end{enumerate}
The fourth class is the most publishable because it says the reduced static
model can choose the wrong reinforcement.

\section{Braess, Braess-Like, and Control-Resonant Reversals}

The literature uses Braess' paradox for cases where adding or increasing
network capacity decreases system performance
\cite{braess2005traffic,schafer2022braess}. In power systems, the performance
loss may be larger flow on an already stressed line, loss of synchrony, or
weaker linear stability \cite{witthaut2012braess,coletta2016braess}. The
present thesis separates three cases.

\begin{definition}[Three reinforcement reversals]
Let \(b_e\) be an upgraded line.
\begin{enumerate}
\item A flow Braess reversal occurs when an upgrade increases loading on a
protected line in the direction of its pre-existing flow.
\item A damping-margin reversal occurs when
\(\partial h_c/\partial b_e<0\) or
\(\partial\zeta_c/\partial b_e<0\) for the protected critical mode.
\item A control-resonant reversal occurs when the line action decreases the
damping of a network-family pole, moves it closer to a condensed controller
pole, and the self-energy term dominates the NEP sensitivity.
\end{enumerate}
\end{definition}

Only the first case is a classical flow Braess statement. The second is a
Braess-like spectral warning. The third is a stronger mechanism that belongs
to the rational self-energy layer:
\begin{equation}
  \frac{ds}{db_e}
  =
  -
  \frac{
  y^*\left(s\,\Sigma_{b_e}(s)+L_{b_e}\right)x
  }{
  y^*T_s(s)x
  }.
  \label{eq:control_resonant_line_sensitivity}
\end{equation}
If the dominant negative contribution comes from the \(s\Sigma_{b_e}(s)\)
channel rather than from \(L_{b_e}\) alone, the reinforcement is not explained
by scalar graph stiffness. It is a network-control resonance.

\section{Internal Validation Gates from the Project}

The project material preserved in \texttt{temp/} contains useful evidence, but
it must be read as gates and diagnostics. The key validated and negative
records are summarized in Table~\ref{tab:internal_validation_gates}.

\begin{table}[ht]
\centering
\caption{Internal validation gates used to calibrate the theory. These are not
universal claims; they define what the later empirical thesis must reproduce
or falsify under registered controls.}
\label{tab:internal_validation_gates}
\begin{tabular}{L{0.27\textwidth}L{0.35\textwidth}L{0.28\textwidth}}
\toprule
\rowcolor{softblue}
\thead{Gate} & \thead{Main record} & \thead{Thesis use}\\
\midrule
Adaptive damping screen &
In the \texttt{ieee39\_rational\_filter} output, 300 surrogate trials gave
median relative error \(4.68\times10^{-3}\) for the diagonal screen,
\(3.21\times10^{-4}\) for the second-order correction, and
\(3.69\times10^{-4}\) for the reduced QEP. &
Use diagonal damping as a fast screen and reduced QEP as a correction, not as
a universal replacement for the full model.\\
\addlinespace
Weak nodes and weak links &
The same surrogate ranked bus 30 as the strongest local damping action
(\(\Delta\zeta=1.12\times10^{-2}\)) and line 25--37 as the strongest positive
line action (\(\Delta\zeta=6.31\times10^{-3}\)); other buses and lines were
negative. &
Weak elements are action- and margin-relative. A centrality list is not enough.\\
\addlinespace
Braess-NEP phase &
The \texttt{phase\_braess\_nep} output tested 460 line-mode perturbations:
348 margin reversals, 232 network-family reversals, and 40 resonant reversals.
The clearest case was line 20--34 with self-energy-dominated
\(\Delta\zeta\) and a reported internal denominator metric of \(0.959\),
which must be recomputed as a magnitude-normalized memory or controller-channel
dominance index. &
Call the effect control-resonant reinforcement risk after NEP validation, not a
generic graph Braess theorem. The missing extraction is the strict subset with
\(d\zeta_{\rm static}/dp\ge0\) and \(d\zeta_{\rm NEP}/dp<0\).\\
\addlinespace
Inertia-control cross term &
The registered full-ANDES gate found 0/13 measurable singleton
GENROU-inertia by PLL-gain cross cases; the strongest cross-dominance index was
\(2.24\times10^{-4}\). &
Keep cross terms in the local theory, but do not use them as the empirical
headline.\\
\bottomrule
\end{tabular}
\end{table}

\FloatBarrier

\section{Measurement Gate for Frequency Security}

Dynamic hosting is also limited by what the protection and estimation layer can
see. The estimator draft in the project compares frequency estimators under
low-inertia IBR conditions and uses the trip-risk exposure
\begin{equation}
  T_{\rm risk}(e_{\rm th})
  =
  \sum_{m=0}^{N-1}\Delta t\,
  \mathbf 1_{\{|e[m]|>e_{\rm th}\}},
  \label{eq:trip_risk_duration}
\end{equation}
where \(e[m]\) is the estimator error relative to the reference frequency.
This metric is included here for one reason: a low RMS estimator can still be
unacceptable if its peak error or latency creates a long protection exposure.
Thus a complete dynamic hosting study should pair spectral margins with
measurement and protection margins. The theoretical hosting set
\eqref{eq:dynamic_hosting_set} can be extended by adding
\(T_{\rm risk}\le T_{\rm max}\) as a protection-aware constraint.

\section{Planning Rule}

The final planning rule is deliberately conservative.

\begin{enumerate}
\item Use static hosting constraints to reject voltage, thermal, and power
quality violations.
\item Use \(\mathcal P_{\rm host}^{(1)}\) to screen SG-to-IBR conversions by
damping headroom.
\item Rank virtual-inertia sites with separate RoCoF and damping-headroom
weights; do not mix the objectives silently.
\item Rank line upgrades by \(G_e^{(2)}\), \(G_e^{(K)}\), and
\(\partial h_k/\partial b_e\), then reject actions that improve graph
stiffness while degrading the protected dynamic margin.
\item If controller states are condensed, verify important candidates with the
NEP sensitivity \eqref{eq:control_resonant_line_sensitivity} or full
eigenvalue recomputation.
\end{enumerate}

This rule gives a defensible version of the practical advice: improve the
graph first when the weak mechanism is graph-mediated, add damping when the
weak mechanism is modal dissipation, add virtual inertia only with the correct
frequency-security objective or co-designed damping, and call a reinforcement
Braess-like only after the measured margin actually worsens.
